\section{Verification architecture}

The canonical release is a compact nested archive rather than a dependency on an online service.  Its outer command is
\begin{verbatim}
python3 verify.py          # quick integrity and retained-certificate audit
python3 verify.py --full   # compile and replay all proof packages
\end{verbatim}
The verification layers are:
\begin{enumerate}
\item SHA-256 of the canonical archive;
\item outer and nested ZIP integrity;
\item package-local SHA-256 manifests;
\item exact retained-certificate audits;
\item recompilation and full finite regeneration where supported;
\item theorem-value consistency checks.
\end{enumerate}
The full clean replay terminates with
\begin{verbatim}
ERDOS506_FULL_REPLAY=PASSED
\end{verbatim}
and requires only Python 3.10 or newer, GNU C++20, Bash, \texttt{sha256sum}, and standard ZIP tools.  No network access is used.  Timeouts, crashes, missing files and missing terminal markers are fail-closed and never interpreted as impossibility.

The canonical release deliberately omits superseded working packages in which timeouts or incomplete case matrices had once been misclassified.  This audit history is important: a computer-assisted proof is only as strong as its completeness accounting and its treatment of failed computation.

